\section*{THE NOTION OF TRUTH IN MATHEMATICS.}
\phantomsection
\addcontentsline{toc}{section}{THE NOTION OF TRUTH IN MATHEMATICS.}

Mathematicians have always been sure that they prove ``truths'' or ``true propositions''; such a conviction can obviously only be sentimental or metaphysical, and it is not by getting on to mathematical ground that it can be justified, nor even given a meaning that does not make it a tautology. The history of the concept of truth in mathematics is the concern therefore of the history of philosophy and not of that of mathematics; but the evolution of this concept has had an undeniable influence on that of mathematics, and because of this we can not let it go by in silence.

Let us note first that it is as rare to see a mathematician in possession of a strong philosophical culture as to see a philosopher who has an extensive knowledge of mathematics; the opinions of mathematicians on topics in philosophy, even when these questions are concerned with their field, are most often opinions received at second or third hand, coming from doubtful sources. But, precisely because of this, it is these average opinions which interest the mathematical historian, at least as much as the original views of thinkers such as Descartes or Leibniz (to mention two who were also mathematicians of the first rank), Plato (who at least kept up with the mathematics of his time), Aristotle or Kant (of whom the same could not be said).

The traditional notion of mathematical truth is that which goes back to the Renaissance. In this concept, there is no great difference between those objects which are the concern of mathematicians and those which are the concern of natural science; both are knowable and man has obtained a grasp on them both through intuition and through reason; there are no grounds to doubt either intuition or reason, which are fallible only if used incorrectly. ``One must have'', says Pascal, ``altogether a wrong spirit to reason badly about principles so great that it is almost impossible for them to escape'' ({[}244{]}, v. XII, p.~9). Descartes, by his fireside, convinced himself that ``there have been only Mathematicians who were able to find some proofs, that is to say some sure and certain reasons'' ({[}85 a{]}, v.VI, p.~19) and that (if his tale is to be believed) well before he had built up a metaphysics in which ``that very same thing'', he says, ``that I previously took as a rule, namely that the objects that we can visualise very clearly and very distinctly are all true, is only ensured because God is or exists, and that he is a perfect being'' ({[}85 a{]}, v. VI, p.~38). If Leibniz objects to Descartes that it cannot be seen how to recognise that an idea is ``clear and distinct'', \(^{23}\) he also considers axioms as obvious and undisputed consequences of the definitions as soon as the expressions are understood. \(^{24}\) It must not be forgotten either that, in the language of that time, mathematics includes many sciences that we do not recognise anymore as such, and sometimes even as far as the art of the engineer; and in the confidence that they inspire, the surprising success of their applications to ``natural philosophy'', to the ``mechanical arts'', to navigation, comes in for a large part.

In this way of looking at things, axioms are no more susceptible to being discussed or put into question than the rules of deduction; at most the choice can be left to each person, according to his preferences, to reason ``in the way of the ancients'' or to give free reign to his intuition. The choice of point of departure is also a matter of individual preference, and one sees appearing numerous ``editions'' of Euclid where the solid logical framework of the Elements becomes a strange travesty; some surveys are given of the infinitesimal calculus and rational mechanics, supposedly deduced from fundamentals which are remarkably badly established; and Spinoza was perhaps giving in good faith his Ethics as being proved in the manner of geometers ``more geometrico demonstrata''. If it is hard to find in the XVIIth century two mathematicians who agree on any matter whatsoever, if the polemics occur daily, endless and acrimonious, the notion of truth remains none the less not in question. ``Only having one truth about each object'', says Descartes, ``whoever finds it knows as much as can be known about it'' ({[}85 a{]}, v. VI, p.~21).

Although no Greek mathematical text of the high period on these questions has survived, it is probable that the point of view of Greek mathematicians on this subject had many more nuances. It is by experience only that the rules of deduction were elaborated to the point of inspiring complete confidence; before they could be considered to be above all discussion, it was necessary to go through mainly fumblings and paralogisms. It would also be to misconstrue the critical spirit of the Greeks, their taste for discussion and sophistry, to imagine that the ``axioms'' themselves that Pascal judged to be the most obvious (and that, according to a legend spread by his sister, he had, with infallible instinct, discovered himself during his childhood) were not the subject of long discussions. In an area that was not that of geometry strictly speaking, the paradoxes of the Eleates have preserved for us some traces of these polemics; and Archimedes when he makes the observation ({[}5 b{]}, v. II, p.~265) that his predecessors made use in many situations of the axiom to which we usually ascribe his name, adds that what is proved by means of this axiom ``has been accepted no less than what has been proved without it'', and that it is sufficient that his own results be accepted on the same basis. Plato, in accordance with his metaphysical views, presents mathematics as a means of access to a ``truth in itself'' and the objects it deals with as having a real existence in the world of ideas; he characterises none the less the mathematical method precisely in a famous extract from the Republic: ``Those that are involved with geometry and arithmetic \ldots{} assume even and odd, three kinds of angles; they treat them as known objects: once assumed, they esteem that they no longer need to give account of them either to themselves or to others, {[}considering it{]} as clear to each one; and starting there, they proceed in sequence, in order to reach by common agreement the goal that their research had suggested'' ({[}250{]}, Book VI, 510 c-e). That which makes up a proof is therefore firstly a point of departure supplying an arbitrary start (even though ``clear to everyone''), and beyond which, he says a bit further on, one does not try to dig; and then, a discussion that follows in sequence a series of intermediate stages; finally, at each step, the consent of the interlocutor guaranteeing the correctness of the reasoning. It must be added that once the axioms are stated, no new appeal to intuition is allowed on principle: Proclus, quoting Géminus, recalls that ``we have learnt from the pioneers of this science themselves, to take no account of conclusions which are only plausible when it is a case of reasoning which must be part of our geometric doctrine'' ({[}153 e{]}, v.I, p.~203).

Thus it is to experience and the cauldron of criticism that is due the elaboration of the rules of mathematical reasoning; and if it is true, as has been argued in a plausible way \([317\ d]\) , that Book VIII of Euclid has kept for us part of the arithmetic of Archytas, it is not surprising to see there the rigidity of the rather pedantic reasoning that does not fail to appear in all mathematical schools where ``rigour'' is discovered or believed to have been discovered. But, once having entered into the practice of mathematicians, it does not seem that these rules of reasoning have ever been doubted until quite recently: if with Aristotle and the Stoics, some of the rules are deduced from others by schemes of reasoning, the primitive rules are always assumed to be evident. In the same way, having gone back to the ``hypotheses'', ``axioms'' and ``postulates'' that appeared to them to supply a solid foundation for the science of their time (such as those for example that they put forward in the first ``Elements'' that tradition ascribes to Hippocrates of Chio, about 450 B. C.), the Greek mathematicians from the classical period seem to have bent all their efforts to the discovery of new results rather than to a critique of these foundations that, at that time, could not have failed to be sterile; and, putting aside all metaphysical preoccupations, it is from this general accord amongst mathematicians on the bases of their science that the text of Plato quoted above bears testimony.

On the other hand Greek mathematicians do not seem to have believed it possible to explain the ``primary notions'' that were their points of departure, straight line, surface, ratio of quantities; if they give ``definitions'' of them, it is obviously a crisis of conscience and without any illusions about their range. It goes without saying that on the other hand, about definitions other than those of ``primary notions'' (definitions often called ``nominal''), the Greek mathematicians and philosophers had perfectly clear ideas. It is in this context that intervenes explicitly, for the first time no doubt, the question of ``existence'' in mathematics. Aristotle does not miss observing that a definition does not imply the existence of the object defined, and that there must be over and above that either a postulate or a proof. No doubt his observation was derived from the practice of mathematicians; in any case Euclid takes care to postulate the existence of a circle, and to prove that of an equilateral triangle, of parallels, of a square, etc. as he introduces them into his arguments ({[}153 e{]}, Book I); these proofs are ``constructions''; in other words, he exhibits, basing himself on the axioms, mathematical objects that he proves satisfy the definitions that he needs to justify.

Thus we see Greek mathematics in the classical era reaching a kind of empirical certainty (whatever the metaphysical bases for it from some philosopher or other might be); if it cannot be conceived that the rules of reasoning can be questioned, the success of Greek science, and the feeling that exists of the inopportunity of a critical revision, play a great part in the confidence that the axioms themselves inspire, a confidence that is rather of the order of that (almost unlimited, as well) that was put during the last century in the principles of theoretical physics. It is in any case what is suggested by the motto of the school ``nihil est in intellectu quod non prius fuerit in sensu'', which is precisely that against which Descartes stands out, as not giving a firm enough basis to that which Descartes hoped to extract from the use of reason.

One must come down to the beginning of the XIXth century to see mathematicians recover from the arrogance of a Descartes (not to say that of a Kant or that of a Hegel, this latter somewhat late, as was fashionable, about the science of his time \(^{25}\) ), to a position as well balanced as that of the Greeks. The first blow to the classical concepts is the establishment of non-Euclidean hyperbolic geometry by Gauss, Lobatschevsky and Bolyai at the beginning of the century. We will not undertake to retrace here in detail the origins of this discovery, the outcome of numerous unfruitful attempts to prove the parallel postulate (see {[}105 a and b{]}). At the time its effect on the principles of mathematics was perhaps not so deep as is sometimes said. It simply forces the abandonment of the pretensions of the previous century to the ``absolute truth'' of Euclidean geometry, and all the more, the Leibnizian point of view that definitions imply axioms; these latter no longer appear as at all ``obvious'', but rather as hypotheses of which it must be determined whether they are adapted to the mathematical representation of the world of the senses. Gauss and Lobatschevsky believe that the debate between the different possible geometries can be determined by experience ({[}206{]}, p.~76). It is also the point of view of Riemann, of whom the famous inaugural Lecture ``On the hypotheses which form the foundation of geometry'' had as aim to furnish a general mathematical framework for the various natural phenomena: ``What remains to be resolved'' he says, ``is the question of knowing to what extent and up to what point these hypotheses are found to be confirmed by experience'' ({[}259 a{]}, p.~284). But that is a problem which visibly has nothing anymore to do with Mathematics; and none of the previous authors seem to put into question that, even if a ``geometry'' does not correspond to an experimental reality, its theorems remain no less ``mathematical truths''. \(^{26}\)

All the same, if it so, it is certainly no longer to an unlimited confidence in classical ``geometric intuition'' that such a conviction is due; the description that Riemann seeks to give of ``multiplicities n times extended'', the object of his work, only relies on ``intuitive'' \(^{27}\) considerations to reach a justification for the introduction of ``local co-ordinates''; from that moment on he apparently seems to feel himself on solid ground, namely that of Analysis. But this latter is grounded in the end on the concept of real number, which remained until then of a very intuitive nature; and progress in the theory of functions was leading in this respect to some very worrying results: with the research of Riemann himself on integration, and more so with the examples of curves without tangents, constructed by Bolzano and Weierstrass, it is the whole of the pathology of mathematics that was beginning. For a century we have seen so many monsters of this species that we are a bit blasé, and the most weird teratological characters must be accumulated in order still to astound us. But the effect produced on the majority of mathematicians of the XIXth century went from disgust to consternation: ``How'' H. Poincaré asks himself, ``can intuition deceive us at this point?'' ({[}251 d{]}, p.~19); and Hermite (not without a spark of humour which the commentators of this famous sentence do not all seem to have perceived) declares that he ``turns away with fear and horror from this lamentable plague of continuous functions that do not have a derivative'' ({[}160{]}, v. II, p.~318). The worst was that these phenomena, so contrary to common sense, could no longer be laid at the door of ideas badly elucidated, as at the time of the ``indivisibles'' (see p.~173), since they survived after the reform of Bolzano, Abel and Cauchy, who had established the foundation of the notion of limit in a manner as rigorous as the theory of ratios (see p.~154). It is thus fully to the gross and incomplete character of our geometric intuition that the account must be laid, and it is reasonable that since then it has remained discredited quite justifiably as a means of proof.

This realisation would inevitably react on classical mathematics, starting with geometry. In whatever respect the axiomatic construction of Euclid had been held, more than one imperfection had been noticed, and that already in antiquity. It was the postulate on parallels that had been the object of the greatest number of criticisms and attempts at proof; but the followers and commentators of Euclid had also attempted to prove other postulates (notably that of the equality of right angles) or recognised the insufficiency of certain definitions, such as those of the straight line or plane. In the XVIth century, Clavius, an editor of the Elements, notes the absence of a postulate guaranteeing the existence of the fourth proportional; for his part, Leibniz remarks that Euclid uses geometric intuition without mentioning it explicitly, for example when he admits (Elements, Book I, prop. 1) that two circles, each of which goes through the centre of the other have a common point ({[}198 b{]}, v. VII, p.~166). Gauss (who himself did not deny himself the use of such topological considerations) draws attention to the role played in Euclidean constructions by the notion of a point (or a straight line) being ``between'' two others, notion that is however not defined ({[}124 a{]}, v. VIII, p.~222). Finally, the use of displacements --- notably in the ``case of equality of triangles'' --- long assumed to be obvious, \(^{28}\) was soon to appear to the critics of the XIXth century as relying also on unstated axioms. One ends up thus, in the period from 1860 to 1885, with different partial revisions of the beginnings of geometry (Helmholtz, Méray, Houël) tending to remedy some of these gaps. But it is only with M. Pasch {[}245{]} that the abandonment of all appeals to intuition is a programme properly formulated and followed with full rigour. The success of his enterprise soon brought him numerous emulators who, principally between 1890 and 1910, gave quite varied statements of the axioms of Euclidean geometry. The most famous of these works were those of Peano, written in his symbolic language {[}246 d{]}, and especially the ``Grundlagen der Geometrie'' of Hilbert {[}163 c{]}, appearing in 1899, a book which, by its lucidity and the depth of its exposition, was to become immediately, with full justification, the charter for modern axiomatics, even to the extent of leading to forgetting its forerunners. It is indeed that, not content with giving there a complete system of axioms for Euclidean geometry, Hilbert classifies these axioms into different groups of different types, and sets himself to determine the exact area of influence of each of these groups of axioms, not only in developing the logical consequences of each of them in isolation, but also in discussing the different ``geometries'' obtained when one omits or modifies certain of these axioms (geometries amongst which those of Lobatschevsky and of Riemann appear only as particular cases \(^{29}\) ); he thus puts clearly in the picture, in an area considered until then as one of those nearest the reality of the senses, the freedom of which the mathematician disposes in his choice of postulates. In spite of the disarray caused in more than one philosopher by these ``metageometries'' with strange properties, the thesis of the ``Grundlagen'' was rapidly adopted almost unanimously by mathematicians; H. Poincaré, although hardly guilty of favouring formalism, recognised in 1902 that the axioms of geometry are conventions, for which the notion of ``truth'', as it is normally understood, has no more meaning ({[}251 c{]}, pp.~66-67). ``Mathematical truth'' resides thus uniquely in logical deduction starting from premises arbitrarily set by axioms. As will be seen later (pp.~35 ff.), the validity of the rules of reasoning used in making these deductions was itself soon to be put in question, bringing about thus a complete reshaping of the basic conceptions of mathematics.

